
\begin{figure*}[p]
    \scriptsize
    \centering
    \begin{tabular}{@{}p{0.31\columnwidth}p{0.31\columnwidth}p{0.31\columnwidth}@{}}
    \toprule
    \textbf{[Stack Exchange (STEM)]}
    
    \textcolor[HTML]{3078BE}{\textbf{What is the difference between minimum and infimum? I have a great confusion about this.}}
    
    ~
    
    The minimum is attained, the infimum isn't necessarily. 

    ~
    
    Example: Let $f(x) = \frac{1}{x}$. Then $f$ has no minimum value on the interval $(0,\infty)$. The minimum is the smallest element in the set. That is $\min\{f(x)\mid x\in (0,\infty)\}$ doesn't exist because there is not smallest number in the set. Another example is the minimum of the set $S = (0,1) = \{x\mid 0<x<1\}$. Here again there isn't a smallest number $\min\{x\mid 0<x<1\}$ doesn't exist.

    ~
    
    The infimum of a set $S$ is defined as the greatest number that is less than or equal to all elements of S (from Wikipedia). The infimum is also sometimes called the greatest lower bound.
    It is a fact that every non empty set (bounded below) of real numbers has an infimum. But, as we saw, not every real set has a minimum.
    So in the example $\inf\{f(x)\mid x\in (0,\infty)\} = 0.$

    ~
    
    Note that the infimum and the minimum can be the same. Consider for example $S = \{1,2,3,\dots\}$. Then the infimum and minimum is both $1$.
    Consider this other example. If $f$ is a continuous function on a closed interval $[a,b]$,  then it is a fact that $f$ attains a minimum over that interval. So here again $\inf\{f(x)\mid x\in [a,b]\} = \min\{f(x)\mid x\in [a,b]\}$.
    
    &
    \textbf{[Stack Exchange (Other)]}
    
    \textcolor[HTML]{3078BE}{\textbf{Was the Millennium Falcon a one-off or was it mass produced?}}
    
    ~
    
    The Millennium Falcon is not unique; it was a model YT-1300f Corellian light freighter. This implies that there are other similar ships, including other variations on the YT-1300 model. You can find a lot of information about the Millennium Falcon in a companion book, Star Wars: The Force Awakens: Incredible Cross-Sections. The model number also appears on-screen at least once: in Star Wars: Rebels episode ``Fighter Flight'' on a poster. Note, however, that the Millennium Falcon has been highly customized; some of the custom work is mentioned in The Force Awakens. It's not clear just how much of the ship is stock and how much is custom.

    ~
    
    If you're willing to go outside Canon into Legends, the model number of the ship, and other similar Corellian light freighters, can be seen in some of the video games, and there is an entire novel called Millennium Falcon that details its history. The ship was serial \#YT-1300 492727ZED, one of many such vessels, but from the beginning the ship was shown to have some ``quirks'' to distinguish it from its counterparts. In Legends, the ship was also destroyed and rebuilt from scratch at least once, so it's arguable as to whether it's still the same ship or not.
    
    &
    \textbf{[wikiHow]}
    
    \textcolor[HTML]{3078BE}{\textbf{How to be a lazy college student?}}

    ~

    In college, unlike in high school, you are expected to take responsibility for your own learning.  It is ultimately up to you to determine how much time and effort you want to put into your schoolwork. Experts suggest that college students today work considerably fewer hours than their counterparts did in 1961.  But there is also evidence that working more hours is not necessarily the best way to be productive. A little laziness may even be good for us.

    ~
    
    \textbf{Part 1 of 3: prioritize what's important}
    \begin{enumerate}[leftmargin=8pt]
        \item Ask yourself what you want from college. There are many reasons why a college degree is desirable, not least that it is still considered highly valuable by many employers. Your own reasons for pursuing a college degree may be more personal, but spend some time working out exactly what these reasons are. Consider why college is important to you. Do you want to learn important skills like critical thinking and problem-solving? Do you want to make friends, network and get to know other people? Do you want to get good grades for further study? Do you want to get a particular job that requires a college degree? Spend some time reflecting on these questions and come up with your own answers. It might even be a good idea to write them down. This will help you to recognize what's important to you so that you can focus on the essential things.
        \item Select your major and your classes carefully. \textcolor[HTML]{808080}{[...]}
    \end{enumerate}

    \\
    \midrule
    \textbf{[Manually Authored -- Chitchat]}
    
    \textcolor[HTML]{3078BE}{\textbf{Tell me an interesting fact about geography.}}

    ~
    
    Here are some interesting facts about geography:
    \begin{itemize}[leftmargin=8pt]
        \item Bhutan is one of the most intriguing places in the world. The Buddhist kingdom is the only carbon-negative country on Earth, with strict environmental laws and cultural norms that preserve its stunning natural landscape. For similar reasons, Bhutan was one of the last countries to introduce television in 1999, and has no traffic lights to this day.
        \item There are currently few territories on Earth that are unclaimed by any sovereign state (terra nullius). One example is Bir Tawil, along the border between Egypt and Sudan, which was created by a discrepancy between borders drawn in 1899 and 1902. The territory remains unclaimed by either country, because it would entail forfeiting the much larger Halaib Triangle, which is claimed by both countries. Bir Tawil has no settled population, but the land is used by Bedouins who roam the area.
        \item The antipode of any spot on Earth \textcolor[HTML]{808080}{[...]}
    \end{itemize} 
    &
    \textbf{[Manually Authored -- Advice]}
    
    \textcolor[HTML]{3078BE}{\textbf{I'm going to NeurIPS for the first time to present the first paper I've ever published. I'm afraid I'll be alone and overwhelmed by everyone. What should I do?}}
    
    ~
    
    Going to an academic conference for the first time may seem like a bewildering experience, but if done right, conferences are a wonderful opportunity to brainstorm with experts from around the world, present your work and get feedback, and make new friends. Here are some tips for improving your conference experience:
    \begin{itemize}[leftmargin=8pt]
    \item Reach out in advance to people who have inspired your work. \textcolor[HTML]{808080}{[...]}
    \item Show curiosity about other people's work. \textcolor[HTML]{808080}{[...]}
    \item Sign up for student volunteer \textcolor[HTML]{808080}{[...]}
    \item Ask your advisor to introduce you \textcolor[HTML]{808080}{[...]}
    \item At the same time, try to socialize outside of your lab \textcolor[HTML]{808080}{[...]}
    \end{itemize}
    
    Ultimately, it is important to remember that every first-conference attendee is probably feeling a bit shy too, and that taking the first step in starting a conversation will almost always be met with a big smile! 

    &
    \textbf{[Manually Authored -- Writing]}
    
    \textcolor[HTML]{3078BE}{\textbf{I am planning to start a book club with some friends. Can you write an email invitation for the same?}}

    ~
    
    Here is an example email invitation for a book club:

    ~

    Subject: Join me for a book club!

    ~
    
    Hi [Friend's Name], how are you?

    ~
    
    I'm starting a book club with some friends, and I'd love to have you join us!

    ~
    
    The book club is a casual and fun way to read and discuss books that we love or want to learn more about, and also make new friends. We plan to meet more-or-less once a month.

    ~
    
    We're thinking of having our first meeting on [Date \& Time]. Could that work for you?

    ~
    
    Would love to see you there!

    ~
    
    Thanks,
    
    [Your Name]
    \\
    \bottomrule
    \end{tabular}
    \caption{\textbf{Training examples from different sources.} \textit{Top row:} examples mined from community Q\&A. \textit{Bottom row:} manually-authored examples. The \textcolor[HTML]{3078BE}{blue} text is the prompt. The gray ellipsis \textcolor[HTML]{808080}{[...]} indicates that the response was trimmed to fit this page, but the actual training example is longer.}
    \label{fig:training_examples}
\end{figure*}
